% !TeX program = xelatex
\documentclass[11pt,a4paper]{report}

% ---------- Encoding, fonts (T5 supports Vietnamese names on the cover) -----
\usepackage{fontspec , float}
\setmainfont{TeX Gyre Termes}
\usepackage{microtype}
\usepackage[english,vietnamese]{babel}

% ---------- Page layout: 2.5 cm margins on all sides ----------
\usepackage[a4paper,margin=2.5cm,headheight=16pt]{geometry}
\usepackage{setspace}
\setstretch{1.0}                 % tài liệu hướng dẫn cho phép tối đa 10 trang
\setlength{\parindent}{0pt}
\setlength{\parskip}{0.3em}
\setlength{\emergencystretch}{2em}

% ---------- Packages ----------
\usepackage{booktabs,tabularx,array,graphicx,float,caption,enumitem,xcolor , subcaption}
\usepackage[most]{tcolorbox}
\usepackage{listings,fancyhdr,titlesec,xurl,hyperref,bookmark}
\usepackage{tikz}
\usetikzlibrary{arrows.meta,positioning}

\definecolor{HCMUSBlue}{HTML}{145DA0}
\definecolor{SoftGray}{HTML}{F3F4F6}

% ---------- Captions: 10 pt ----------
\captionsetup{font=footnotesize,labelfont=bf,justification=centering}

\setlist[itemize]{leftmargin=1.5em,itemsep=0pt,topsep=0.3em}
\setlist[enumerate]{itemsep=0pt,topsep=0.3em}
% ---------- Syntax highlighting for code listings ----------
\definecolor{CodeKeyword}{HTML}{1F4FBF}   % từ khóa C
\definecolor{CodeType}{HTML}{267F99}      % kiểu dữ liệu
\definecolor{CodeComment}{HTML}{2E7D32}   % chú thích
\definecolor{CodeString}{HTML}{B3261E}    % chuỗi và ký tự
\definecolor{CodeNumber}{HTML}{B45F06}    % số
\definecolor{CodePrep}{HTML}{7B1FA2}      % #include, #define
\definecolor{CodeSyscall}{HTML}{C2185B}   % lời gọi hệ thống
\definecolor{CodeLib}{HTML}{00796B}       % hàm thư viện
\lstdefinestyle{ccode}{
  language=C,
  basicstyle=\ttfamily\scriptsize,
  keywordstyle=\color{CodeKeyword}\bfseries,
  keywordstyle=[2]\color{CodeType},
  commentstyle=\color{CodeComment}\itshape,
  stringstyle=\color{CodeString},
  numberstyle=\color{CodeNumber},
  directivestyle=\color{CodePrep}\bfseries,
  emphstyle=\color{CodeSyscall}\bfseries,
  emphstyle=[2]\color{CodeLib},
  morekeywords=[2]{uint,uint64,uchar,ushort},
  emph={interpose,sbrk,open,exec,exit,fork,wait,read,write},
  emph=[2]{argstr,argint,strncmp,safestrcpy,memset,memcmp,strlen,printf,fprintf},
  breaklines=true,frame=single,rulecolor=\color{black!25},
  columns=fullflexible,backgroundcolor=\color{SoftGray},showstringspaces=false}
\lstdefinestyle{terminal}{
  language={},keywordstyle={},emph={},emph=[2]{},commentstyle={},stringstyle={},
  basicstyle=\ttfamily\scriptsize,
  moredelim=[l][\color{HCMUSBlue}\bfseries]{\$},
  breaklines=true,frame=single,rulecolor=\color{black!25},
  columns=fullflexible,backgroundcolor=\color{SoftGray},showstringspaces=false}
\lstset{style=ccode,aboveskip=3pt,belowskip=0pt}

  % ---------- Header / footer ----------
\pagestyle{fancy}
\fancyhf{}
\fancyhead[R]{\small Project 2: Griductive Solver}
\fancyfoot[C]{\thepage}
\renewcommand{\headrulewidth}{0.4pt}

% ---------- Headings: 16 pt chapter, 14 pt section, 12 pt subsection (all bold) ----------
\titleclass{\chapter}{straight}
\titleformat{\chapter}[display]{\bfseries\fontsize{16}{20}\selectfont}{\filright\MakeUppercase{Chương \thechapter}}{0.5ex}{\titlerule\vspace{1ex}\filright}
\titlespacing*{\chapter}{0pt}{10pt}{6pt}
\titleformat{\section}{\bfseries\fontsize{14}{17}\selectfont}{\thesection}{0.7em}{}
\titlespacing*{\section}{0pt}{10pt}{4pt}
\titleformat{\subsection}{\bfseries\fontsize{12}{15}\selectfont}{\thesubsection}{0.7em}{}
\titlespacing*{\subsection}{0pt}{6pt}{2pt}
\setcounter{tocdepth}{1}

\hypersetup{colorlinks=true,linkcolor=HCMUSBlue,urlcolor=HCMUSBlue,
  pdftitle={Project 2: Griductive Solver},pdfauthor={Nhóm sinh viên lớp 24CNTN}}
\renewcommand{\contentsname}{MỤC LỤC}
\newcolumntype{Y}{>{\raggedright\arraybackslash}X}

% ---------- Edit only these fields to update the cover and contribution table ----------
\newcommand{\ClassName}{24CNTN}
\newcommand{\GroupID}{[Mã nhóm]}
\newcommand{\InstructorName}{[Tên giảng viên]}
\newcommand{\MemberOne}{Huỳnh Bảo Khang }
\newcommand{\IDOne}{24120507}
\newcommand{\MemberTwo}{Phạm Minh Đạt }
\newcommand{\IDTwo}{24120174}
\newcommand{\MemberThree}{Huỳnh Nguyên Phát }
\newcommand{\IDThree}{24120005}
\newcommand{\MemberFour}{Trần Thành Nhân }
\newcommand{\IDFour}{24120402}
\newcommand{\MemberFive}{Nguyễn Phạm Xuân Cương }
\newcommand{\IDFive}{24120499}
\newcommand{\SubmissionDate}{[Ngày nộp]}
\newcommand{\cmd}[1]{\texttt{#1}}    % lệnh: make qemu
\newcommand{\file}[1]{\texttt{#1}}   % tên file: user/sleep.c
\newcommand{\term}[1]{\emph{#1}}     % thuật ngữ lần đầu

% Đánh dấu đỏ cho mọi thứ vẫn cần điền hoặc kiểm tra trước khi nộp.
% Tìm kiếm "TODO" trong tệp và xóa tất cả.
\newcommand{\TODO}[1]{\textcolor{red}{\textbf{[TODO: #1]}}}

% ---------- Logo: use forward slashes and match the filename exactly ----------
\newcommand{\LogoFile}{assets/hcmus-logo.png}

\newcommand{\SchoolLogo}{%
    \IfFileExists{\LogoFile}{%
        \includegraphics[width=4cm,height=3.3cm,keepaspectratio]{\LogoFile}%
    }{%
        \fbox{%
            \begin{minipage}[c][3cm][c]{3.8cm}
                \centering
                \small LOGO HCMUS\par
                \vspace{0.2cm}
                \footnotesize Tải logo trường lên\par
                theo đường dẫn trong\par
                \texttt{\string\LogoFile}.
            \end{minipage}%
        }%
    }%
}

% ---------- Helper boxes ----------
\newcommand{\FillIn}[2][0.65cm]{%
\begin{tcolorbox}[colback=SoftGray,colframe=black!25,boxrule=0.4pt,
  arc=0pt,left=7pt,right=7pt,top=6pt,bottom=6pt]
{\small\setstretch{1.15}\itshape #2\par}\vspace{#1}
\end{tcolorbox}}

% Đặt ảnh chụp màn hình vào figures/. Tệp bị thiếu sẽ được thay thế bằng một hộp có thể chỉnh sửa.
% \Shot[width]{file name}{caption}
\newcommand{\Shot}[3][0.95\linewidth]{%
\begin{figure}[H]\centering
\IfFileExists{figures/#2}{\includegraphics[width=#1]{figures/#2}}{\FillIn[0.3cm]{Chèn ảnh chụp màn hình \texttt{figures/#2}: #3}}
\caption{#3}
\end{figure}}

% Hai ảnh chụp màn hình cạnh nhau trong một hình: \ShotPair{left file}{right file}{caption}
\newcommand{\ShotPair}[3]{%
\begin{figure}[H]\centering
\begin{minipage}[t]{0.49\linewidth}\centering
\IfFileExists{figures/#1}{\includegraphics[width=\linewidth]{figures/#1}}{\FillIn[0cm]{Chèn \texttt{figures/#1}}}
\end{minipage}\hfill
\begin{minipage}[t]{0.49\linewidth}\centering
\IfFileExists{figures/#2}{\includegraphics[width=\linewidth]{figures/#2}}{\FillIn[0cm]{Chèn \texttt{figures/#2}}}
\end{minipage}
\caption{#3}
\end{figure}}

\begin{document}
\begin{titlepage}

\centering
{\Large\bfseries ĐẠI HỌC QUỐC GIA THÀNH PHỐ HỒ CHÍ MINH}\\[0.15cm]
{\large\bfseries TRƯỜNG ĐẠI HỌC KHOA HỌC TỰ NHIÊN}\\[0.25cm]
{\normalsize\bfseries KHOA CÔNG NGHỆ THÔNG TIN}\\[0.35cm]
{\normalsize Lớp: \ClassName}\\[0.65cm]

\SchoolLogo
\par\vspace{0.7cm}

\rule{\textwidth}{1.5pt}\\[0.55cm]
{\fontsize{24}{30}\selectfont\bfseries BÁO CÁO THỰC HÀNH}\\[0.4cm]
{\Large\bfseries Project 2: Griductive Solver}\\[0.55cm]
\rule{\textwidth}{1.5pt}\\[0.45cm]
{\large\bfseries Môn học: Cơ sở Trí tuệ Nhân tạo}\par
\vspace{0.25cm}

\vspace{0.55cm}

\begin{center}
\renewcommand{\arraystretch}{1.3}
\begin{tabular}[t]{@{}l@{}}
\textbf{Giảng viên hướng dẫn:}\\[0.15cm]
GS.TS Lê Hoài Bắc     \\      
TS. Nguyễn Ngọc Đức  
\end{tabular}
\hspace{2.5cm}
\begin{tabular}[t]{@{}l@{}}
\textbf{Thành viên nhóm:}\\[0.15cm]
\MemberOne (\IDOne) \\
\MemberTwo (\IDTwo)\\
\MemberThree (\IDThree) \\
\MemberFour (\IDFour)\\
\MemberFive (\IDFive) \\

\end{tabular}
\end{center}
\end{titlepage}
\tableofcontents
\clearpage

% =====================================================
%  BÁO CÁO CHÍNH
% =====================================================
\pagenumbering{arabic}
\parindent=1.5em \leftskip=0pt \rightskip=0pt \parfillskip=0pt plus 1fil


\chapter{Giới thiệu và Phân công công việc}

\section{Mục tiêu và Phạm vi}
\noindent\textbf{Mục tiêu.}
Griductive là trò chơi suy luận trên lưới, trong đó mỗi nhân vật bí mật là
\texttt{CRIMINAL} hoặc \texttt{INNOCENT}. Nhóm xây dựng bộ giải không đoán mò:
một phán quyết chỉ được đưa ra khi trạng thái đó đúng trong mọi mô hình thỏa các
clue công khai. Hệ thống sử dụng logic mệnh đề, mã hóa CNF và bộ giải SAT DPLL
tự cài đặt.

\noindent\textbf{Phạm vi.}
Báo cáo mô tả kiến trúc Engine--Encoder--DPLL--Agent--GUI, thuật toán DPLL,
quy trình kiểm thử đối chiếu brute-force và thực nghiệm trên cùng một corpus
gồm các lưới $3\times3$, $4\times4$ và $5\times5$.

\section{Phân công công việc trong nhóm}

\begin{table}[H]
\centering\footnotesize\renewcommand{\arraystretch}{1.15}
\begin{tabularx}{\textwidth}{@{}p{2.1cm}>{\raggedright\arraybackslash}p{3.3cm}Y>{\centering\arraybackslash}p{1.6cm}@{}}
\toprule
\textbf{MSSV} & \textbf{Họ và tên} & \textbf{Nhiệm vụ} & \textbf{Hoàn thành}\\
\midrule
\IDOne & \MemberOne & Logic Agent, entailment, uniqueness, deduction loop và tích hợp & 100\%\\
\IDTwo & \MemberTwo & Game Engine, clue model, semantic evaluator và bộ puzzle & 100\%\\
\IDThree & \MemberThree & DPLL, assumptions, thống kê và experiments & 100\%\\
\IDFour & \MemberFour & CNF encoder, suy dẫn mã hóa và tổng hợp báo cáo & 100\%\\
\IDFive & \MemberFive & GUI, kiến trúc Engine/Agent và tương tác người dùng & 100\%\\
\bottomrule
\end{tabularx}
\caption{Phân công nhiệm vụ và mức hoàn thành do nhóm ghi nhận.}
\end{table}

\chapter{Báo cáo các nhiệm vụ}

\section{Kiến trúc và ranh giới dữ liệu}
Game Engine giữ puzzle đầy đủ, nghiệm ẩn và các clue chưa mở. Agent chỉ nhận
\texttt{public\_state} gồm layout, trạng thái đã biết và clue đã công khai.
Encoder chuyển trạng thái này thành CNF; DPLL giải các truy vấn SAT; Agent dùng
hai truy vấn đối nghịch để phân loại từng ô:
\begin{itemize}
  \item $KB \land \neg C_i$ UNSAT $\Rightarrow C_i$ là \texttt{CRIMINAL};
  \item $KB \land C_i$ UNSAT $\Rightarrow C_i$ là \texttt{INNOCENT};
  \item cả hai SAT $\Rightarrow$ trạng thái \texttt{UNKNOWN}.
\end{itemize}
Ranh giới này bảo đảm Agent không đọc đáp án bí mật và không thể hợp thức hóa
một phán quyết chỉ vì nó tình cờ trùng nghiệm authored.

\section{Thuật toán DPLL}

\subsection{Giao diện và biểu diễn}
Literal được biểu diễn bằng số nguyên có dấu; biến $i$ tương ứng ô thứ $i$ theo
thứ tự hàng--cột. Công thức CNF là danh sách các clause. Giao diện chính là
\texttt{solve(clauses, n\_vars, assumptions=())}; kết quả gồm cờ SAT, model đầy
đủ và thống kê \texttt{decisions}, \texttt{propagations}, \texttt{backtracks},
\texttt{runtime}.

\subsection{Quy trình giải}
\begin{enumerate}
  \item Nạp assumptions và phát hiện cặp literal mâu thuẫn.
  \item Thực hiện unit propagation đến điểm cố định. Clause chưa thỏa không còn
        literal tự do tạo xung đột.
  \item Nếu mọi clause đã thỏa, điền các biến chưa gán bằng giá trị mặc định để
        trả model đầy đủ.
  \item Chọn biến chưa gán có chỉ số nhỏ nhất, thử \texttt{True} trước rồi
        \texttt{False}; quay lui khi một nhánh xung đột.
  \item Ở biến thể so sánh, pure-literal elimination được chạy sau unit
        propagation trên các clause chưa thỏa.
\end{enumerate}

\begin{lstlisting}[style=terminal]
DPLL(F, A):
  A <- UNIT_PROPAGATE(F, A)
  if conflict: return UNSAT
  if every clause is satisfied: return complete_model(A)
  v <- smallest unassigned variable
  for value in [True, False]:
      result <- DPLL(F, A union {v=value})
      if result is SAT: return result
  return UNSAT
\end{lstlisting}

Unit propagation và pure-literal elimination bảo toàn tính khả thỏa. Khi phân
nhánh, hai trường hợp $v=\texttt{True}$ và $v=\texttt{False}$ phủ toàn bộ model;
vì vậy chỉ kết luận UNSAT sau khi cả hai nhánh thất bại. Việc chọn biến nhỏ nhất
giúp kết quả tái lập giữa các lần chạy.

\section{Thiết kế thực nghiệm}
Corpus cố định gồm 4 bài $3\times3$, 4 bài $4\times4$ và 2 bài $5\times5$ làm
stress test. Bốn benchmark bổ sung nằm trong \texttt{experiments/cases.py}; mỗi
bài được kiểm nghiệm duy nhất và phải giải trọn vẹn. Hai biến thể DPLL cơ bản và
DPLL có pure-literal chạy trên đúng cùng corpus. Runner giữ lại mọi trạng thái
\texttt{COMPLETE}, \texttt{STUCK}, \texttt{INCONSISTENT}, \texttt{TIMEOUT} hoặc
\texttt{ERROR} thay vì loại bỏ ca thất bại.

Với mỗi lượt chạy, hệ thống ghi số biến chính/phụ, clause, SAT call, decisions,
propagations, backtracks, số bước reveal và runtime vào
\texttt{experiments/results.csv}. Encoder dùng biến phụ khi cần cho các ràng
buộc mở rộng; trên corpus hiện tại, $n_{aux}$ nằm trong khoảng 0--40.

\chapter{Kiểm thử và Xác minh}

\section{Môi trường và Quy trình kiểm thử}
\begin{table}[H]
\centering\footnotesize
\begin{tabularx}{\textwidth}{@{}p{4.4cm}Y@{}}
\toprule
\textbf{Thành phần} & \textbf{Giá trị}\\
\midrule
Hệ điều hành & Ubuntu 24.04 trên WSL2\\
CPU & 13th Gen Intel Core i5-13420H, 12 logical CPUs\\
Python & 3.12.3\\
Chiến lược chọn biến & Biến chưa gán nhỏ nhất; True trước, False sau\\
Timeout & 30 giây/puzzle, kiểm giữa các deduction step\\
\bottomrule
\end{tabularx}
\caption{Môi trường thực nghiệm.}
\end{table}

Trước khi đo, solver được đối chiếu brute-force trên 200 CNF ngẫu nhiên 1--5
biến; test bao phủ SAT/UNSAT, assumptions, model đầy đủ và uniqueness. Bốn
benchmark đều \texttt{COMPLETE}; regression suite đạt \textbf{123/123 test}.

\section{Kết quả}
\begin{table}[H]
\centering\scriptsize\setlength{\tabcolsep}{3.2pt}
\resizebox{\textwidth}{!}{%
\begin{tabular}{@{}lrrrrrrrrr@{}}
\toprule
\textbf{Puzzle} & $N$ & \textbf{Var} & \textbf{Clause} & \textbf{SAT} &
\textbf{Decision} & \textbf{Propagation} & \textbf{Backtrack} &
\textbf{Step} & \textbf{ms}\\
\midrule
village-3 & 3 & 9 & 27 & 63 & 30 & 363 & 0 & 7 & 3.798\\
village-4 & 4 & 16 & 185 & 224 & 43 & 2163 & 0 & 14 & 17.359\\
village-5 & 5 & 25 & 304 & 575 & 1472 & 7986 & 0 & 23 & 221.490\\
village-5-mist & 5 & 25 & 218 & 575 & 1409 & 8547 & 5 & 23 & 187.484\\
village-3-links & 3 & 9 & 121 & 63 & 81 & 1606 & 6 & 7 & 18.329\\
village-4-links & 4 & 16 & 519 & 224 & 588 & 12217 & 133 & 14 & 540.917\\
benchmark-3-river & 3 & 9 & 16 & 63 & 0 & 350 & 0 & 7 & 2.764\\
benchmark-3-market & 3 & 9 & 16 & 63 & 0 & 350 & 0 & 7 & 4.448\\
benchmark-4-tower & 4 & 16 & 29 & 224 & 0 & 1827 & 0 & 14 & 15.775\\
benchmark-4-bridge & 4 & 16 & 28 & 224 & 0 & 1827 & 0 & 14 & 13.958\\
\bottomrule
\end{tabular}}
\caption{Kết quả DPLL cơ bản trên corpus cố định.}
\end{table}

\begin{table}[H]
\centering\footnotesize
\begin{tabular}{@{}lrrrrrr@{}}
\toprule
\textbf{Biến thể} & \textbf{Hoàn tất} & \textbf{SAT call} & \textbf{Decision} &
\textbf{Propagation} & \textbf{Backtrack} & \textbf{Tổng ms}\\
\midrule
DPLL & 10/10 & 2298 & 3623 & 37236 & 144 & 1026.322\\
DPLL + pure-literal & 10/10 & 2298 & 3623 & 37252 & 144 & 1482.804\\
\bottomrule
\end{tabular}
\caption{So sánh hai biến thể trên cùng corpus.}
\end{table}

\begin{figure}[H]
\centering
\begin{tikzpicture}[x=1cm,y=1cm]
\draw[black!35] (0,-0.05) -- (0,1.45);
\fill[HCMUSBlue] (0,0.80) rectangle (4.15,1.25);
\fill[orange!85!red] (0,0.10) rectangle (5.99,0.55);
\node[anchor=east,font=\footnotesize] at (-0.12,1.02) {DPLL};
\node[anchor=east,font=\footnotesize] at (-0.12,0.32) {DPLL + pure};
\node[anchor=west,font=\footnotesize] at (4.25,1.02) {1026.322 ms};
\node[anchor=west,font=\footnotesize] at (6.09,0.32) {1482.804 ms};
\end{tikzpicture}
\caption{Tổng runtime của hai biến thể trong một lần chạy WSL2.}
\end{figure}

Cả hai biến thể hoàn tất 10/10 puzzle. Bài \texttt{village-4-links} và hai bài
$5\times5$ chiếm phần lớn decisions và runtime. Pure-literal không giảm
decisions trên corpus này;
việc quét polarity chỉ tạo thêm overhead. Các benchmark chuỗi có 0 decision vì
clue quan hệ/đếm kết hợp với trạng thái đã biết đủ để unit propagation suy ra
toàn bộ đáp án.

\chapter{Thảo luận và Tài liệu tham khảo}

\section{Khó khăn, Lỗi gặp phải và Cách khắc phục}
\begin{itemize}
  \item \textbf{Thiếu dữ liệu đúng kích thước:} repo ban đầu chỉ có một bài
        $3\times3$ và một bài $4\times4$. Nhóm bổ sung bốn benchmark chỉ dùng cho
        experiments và kiểm tra tính duy nhất trước khi đo.
  \item \textbf{Độ tin cậy của DPLL:} solver được đối chiếu với brute-force trên
        CNF ngẫu nhiên thay vì chỉ kiểm các ví dụ thuận lợi.
  \item \textbf{So sánh công bằng:} mọi biến thể dùng cùng corpus và giữ lại
        trạng thái lỗi/timeout trong CSV.
  \item \textbf{Khả năng tái lập:} chọn biến deterministic, lưu đầy đủ metrics
        và sinh bảng báo cáo từ kết quả chạy thực tế.
\end{itemize}

Các hạn chế còn lại gồm: mã hóa đếm bằng tổ hợp trực tiếp có thể làm số clause
tăng mạnh; benchmark bổ sung dễ hơn hồ sơ authored; runtime mới đo một lần và
chịu nhiễu WSL2; timeout chưa ngắt được một SAT call đang chạy. Hướng phát triển
là dùng sequential counter/cardinality network, chạy lặp để báo median và dùng
process riêng để thực thi deadline cứng.

\section{Đánh giá kết quả}
Hệ thống tích hợp Game Engine, CNF Encoder, bộ giải DPLL và Logic Agent theo
ranh giới dữ liệu đã thiết kế. Bộ giải đáp ứng giao diện \texttt{solve},
assumptions, model đầy đủ và bốn thống kê theo yêu cầu. Runner experiments tạo
bảng cho 4 bài $3\times3$ và 4 bài $4\times4$, đồng thời giữ 2 bài
$5\times5$ làm stress test. Toàn bộ 123 test vượt qua và 20/20 lượt chạy của hai
biến thể đạt \texttt{COMPLETE}. Kết quả chưa đủ để khẳng định pure-literal luôn
kém hơn; nhóm chỉ kết luận heuristic này chưa có lợi trên corpus hiện tại.

\section{Tái lập kết quả}
\begin{lstlisting}[style=terminal]
cd GriductiveSolver
.venv/bin/python -m pytest -q
.venv/bin/python -m experiments.run
\end{lstlisting}
Runner sinh \texttt{experiments/results.csv}, \texttt{summary.md} và
\texttt{runtime.svg}. Báo cáo này dùng đúng số liệu trong CSV tại thời điểm
25/11/2026.

\section{Tài liệu tham khảo và Khai báo sử dụng công cụ AI}
\begingroup\footnotesize
\begin{enumerate}[leftmargin=1.5em,itemsep=1pt,topsep=2pt]
\footnotesize
\item Đề bài CSC14003 -- Project 2: \textit{Griductive Solver}, các mục
      4.3--4.6 về DPLL, Logic Agent, Experiments và báo cáo.
\item M. Davis, G. Logemann, D. Loveland, ``A Machine Program for
      Theorem-Proving'', \textit{Communications of the ACM}, 1962.
\item Mã nguồn nội bộ của nhóm: \texttt{game/dpll.py},
      \texttt{game/stub\_sat.py}, \texttt{game/agent.py} và
      \texttt{experiments/run.py}, truy xuất ngày 25/11/2026.
\item OpenAI Codex/ChatGPT được sử dụng để hỗ trợ triển khai solver, tạo test,
      tổ chức experiments và trình bày báo cáo. Nhóm đã kiểm chứng bằng
      brute-force, test tự động và CSV; không có số liệu thực nghiệm nào được
      điền thủ công.
\end{enumerate}
\endgroup
\end{document}
